\documentclass[11pt]{article}
\usepackage{fontspec}
\setmainfont{Times New Roman}
\setsansfont{Arial}
\setmonofont{Consolas}
\usepackage{amsmath,amssymb,mathtools}
\usepackage{geometry}
\usepackage{microtype}
\geometry{a4paper,margin=2.05cm}
\setlength{\parindent}{1.1em}
\setlength{\parskip}{0pt}
\newcommand{\foreign}[1]{#1}
\emergencystretch=220em
\allowdisplaybreaks
\sloppy
\hbadness=10000
\vbadness=10000
\hfuzz=5pt
\vfuzz=12pt
\raggedbottom
\begin{document}

\section*{39. Hiperkompleksne sistemy v svojih odnosah s komutativnoju algebraju i teorijeju čisel}

\begin{center}
\foreign{\emph{Verhandl. Intern. Math.-Kongreß Zürich 1 (1932), s. 189--194}}
\end{center}

\begin{center}
\textsc{Emmy Noether, Göttingen}
\end{center}

\paragraph{1.}
Teorija hiperkompleksnyh sistemov, to jest algeber, v poslednjih godah dostala silny vzmah; ale tek v najnovějšem vremeni stalo jasno značenje toj teorije za komutativne voprosy. O tom značenju nekomutativnogo za komutativno hču dnes govoriti: imenno hču podrobno prěslěditi jego na dvuh klasičnyh voprosah, idućih nazad k Gaussu, na teoremu o glavnom rodu i na těsno s nim svęzanoj normnoj teoremě. Te voprosy v tečenju vrěmena znovu i znovu měnjali svoju formulaciju: pri Gaussu pojavjajut se kako zaključok jego teorije kvadratnyh form; potom igrajut suščstvenu rolu v harakterizaciji relativno cikličnyh i abelovyh čislovyh polj posrědstvom teorije klasovyh polj; nakonec možno jih izraziti kako teoremy o avtomorfizmah i o razloženju algeber, a ta poslednja formulacija jednočasno daje prěnos teoremov na libovoljne relativno Galoisove čislove polja.

S toj skicoju, ktoru pozdněje razširju, hču jednočasno objasniti \emph{princip} priměnjenja nekomutativnogo k komutativnomu: \emph{posrědstvom teorije algeber iščut se invariantne i proste formulacije za znane fakty o kvadratnyh formah ili cikličnyh poljah, to jest take formulacije, ktore zavisęt samo od strukturnyh svojstv algeber. Ako sut te invariantne formulacije uže dokazane -- a imenno tako jest v vyše navedenyh priměrah --, togda s tym samym avtomatično dobivaje se prěnos tyh faktov na libovoljne Galoisove polja.}

\paragraph{2.}
Prěd podrobnym izloženjem hču dati ještě obči pregled različnyh metodov i daljnyh rezultatov. Najprvo treba zamětiti, že glavna težkost leži v dobyvanju formulacije za obče Galoisove polja; bez hiperkompleksnoj metody tut obće ne bylo izhodiščnoj točky. V navedennyh priměrah osnovny prěhod k nekomutativnomu dobivaje se črez \emph{jednočasno razsmotrenje polja i grupy} posrědstvom ``skrěženogo produkta'' i jego multiplikativnyh konstant, ``faktornyh sistemov'' (sr. 3). Tako se dobivaje prosta normalna algebra nad osnovnym poljem, i vsaka taka algebra v suščnosti može byti porođena takim sposobom. Take skrěžene produkty prvi razsmatrjal Dikson,\footnote{Sr. jego knigu \foreign{\emph{Algebren und ihre Zahlentheorie}}, Zürich 1927, § 34.} medžutym kak teorija faktornyh sistemov razvila se pri Špajzeru, Šuru i R. Braueru\footnote{Sr. napr. R. Brauer, \foreign{\emph{Untersuchungen über die arithmetischen Eigenschaften von Gruppen linearer Substitutionen}}, i tam v prim. 2 navedenu literaturu, Math. Z. 28 (1928).} iz sovsěm drugoj izhodiščnoj točky, imenno iz voprosa absolutno nerazložimyh predstavjenj. Tek slitje oboju teorij dalo dostatočno prostu i široku konstrukciju, s ktoroju možno bylo pristupiti k komutativnym voprosam.\footnote{Tu konstrukciju ja najprvo razvila v lekciji zimoju 1929/30; ona jest vozproizvedena v gl. 2 raboty H. Hasse, \foreign{\emph{Theory of cyclic algebras}}, Trans. Amer. Math. Soc. 34 (1932). O cěloj oblasti, razsmatrjanoj v dokladě, orientuje pregled M. Deuringa o hiperkompleksnyh čislah i čislovo-teoretičnyh priměnjenjah, ktory ima iziti v seriji \foreign{\emph{Ergebnisse der Mathematik}}.}

Pri tom jednočasno nastavajut nove i prozračne dokazy za znane fakty. Hču tut ukazati na hiperkompleksny dokaz zakona vzajemnosti za ciklične polja, ktory dal H. Hasse i ktory skoro pojavi se v \foreign{Math. Ann.},\footnote{H. Hasse, \foreign{\emph{Die Struktur der R. Brauerschen Algebrenklassengruppe über einem algebraischen Zahlkörper (insbesondere Normenrestsymbol und Reziprozitätsgesetz)}}. Math. Ann. 107 (1932/33).} posrědstvom invariantnoj formulacije svojego normno-ostatkovogo simbola, osnovanoj na teoriji skrěženogo produkta. Dalje ukazujem na hiperkompleksno obosnovanje lokalnoj teorije klasovyh polj, takodže na toj samoj osnově, ktore nedavno dal C. Ševalle; pri tom však bylo treba razviti ještě nove algebraične teoremy o faktornyh sistemah.\footnote{C. Chevalley, \foreign{\emph{Sur la théorie du symbole de restes normiques}}; ima iziti v J. f. Math. 169.} Jednočasno musju však ograničitelno zamětiti, že metoda skrěženyh produktov sama po sebě, po vsemu vidu, \emph{ne daje polnoj teorije Galoisovyh čislovyh polj}. To slěduje iz novyh, ještě nepublikovanyh rezultatov Artina, ktore prodolžajut vyše navedeni dokaz Hasse v duhu označenogo principa, ale dajut samo ravnosti čislovosti vměsto polnyh teoremov o izomorfizmu.

Metody, ktore dajut polny izomorfizm -- pravda, operatorny izomorfizm -- imajut se teper uže v samoj algebrě. Ide o prodolženje pristupov A. Špajzera,\footnote{A. Speiser, \foreign{\emph{Gruppendeterminante und Körperdiskriminante}}. Math. Ann. 77 (1916).} imenno o ponimanje Galoisovogo polja kako ``Galoisovogo modula'', to jest kako modula nad osnovnym poljem, ktory dopušča substitucije Galoisovej grupy kako operatory. I obstaja operatorny izomorfizm medžu poljem i grupnym kolcom (grupnoju algebraju) v tom smyslu, že medžu elementami imaje město vzajemno jednoznačno odpovědanje tak, že linearne formy nad osnovnym poljem odpovědajut jedne drugim, i že substitucijam Galoisovej grupy v polju odpovědaje množenje v grupnom kolcu. Tu teoremu, formulovanu mnoju,\footnote{E. Noether, \foreign{\emph{Normalbasis bei Körpern ohne höhere Verzweigung}}, teorema 3, J. f. Math. 167 (1932) (dokaz soderži prazdninu).} dokazal M. Deuring,\footnote{M. Deuring, \foreign{\emph{Galoissche Theorie und Darstellungstheorie}}. Math. Ann. 107 (1932).} i na njej osnoval dokaz Galoisovej teorije, gde operatorny izomorfizm realizuje odpovědanje medžu grupoju i poljem. Daljne strukturne teoremy -- takodže Deuringove -- idut paralelno formalnym faktam o Artinovyh \(L\)-rędah i dajut strukturny pristup k Artinovym provodnikam. Ti Artinovi \(L\)-rędy i provodniki,\footnote{E. Artin, \foreign{\emph{Über eine neue Art von \(L\)-Reihen}}, Math. Sem. Hamburg 3 (1924). \foreign{\emph{Zur Theorie der \(L\)-Reihen mit allgemeinen Gruppencharakteren}}, tam samo 8 (1931). \foreign{\emph{Die gruppentheoretische Struktur der Diskriminanten algebraischer Zahlkörper}}, J. f. M. 164 (1931).} utvorjene s občimi grupnymi karakterami, pored Špajzera\textsuperscript{6)} sostavljajut prvu svęz medžu teorijeju čisel i teorijeju predstavjenj, prvi prodor poza abelove polja. Oni dali cělomu razvitju silny impuls; osoblivo teorija Galoisovyh modulov orientovala se po njih.

\paragraph{3.}
Teper hču podrobno prěslěditi problemy postavjene na početku: normny teorem i teorem o glavnom rodu. Najprvo \emph{definicija skrěženogo produkta}. Nehaj \(K/k\) bude Galoisovo polje \(n\)-togo stepena, \(\mathfrak G\) jego grupa. Skrěženy produkt označa jednočasno vstavjenje \(K\) i \(\mathfrak G\) v algebru \(A\) tak, že avtomorfizmy od \(K\) stavajut vnutrišnimi. Nehaj \(u_{S_1},\ldots,u_{S_n}\) označajut symboly, odpovědajuče \(n\) elementam grupy; togda najprvo postavja se \(A\) kako modul linearnyh form ranga \(n\) nad \(K\):
\[
\begin{gathered}
(1)\qquad A=u_{S_1}K+\cdots+u_{S_n}K\\
\text{(t. j. \(A\) sostoji iz vsih linearnyh form}\\
\text{\(u_{S_1}a_1+\cdots+u_{S_n}a_n\) s libovoljnymi \(a_i\) v \(K\)).}
\end{gathered}
\]
Črez trebovanje vnutrišnjego avtomorfizma -- porođenogo elementami \(u_S\), občneje \(u_SK^\ast\)\footnote{\(K^\ast\) nastavaje iz \(K\) izostavjenjem nuly; ta oznaka upotrěbljaje se obče.} -- \(A\) stavaje kolcom, znači algebraju ranga \(n^2\) nad \(k\). To trebovanje imenno izraža se črez
\[
(2)\qquad u_S^{-1}zu_S=z^S,\qquad \text{ili}\qquad zu_S=u_Sz^S
\]
za vsako \(z\in K\),\footnote{\(z^S\) označa, kako obyčno, element, nastajuči iz \(z\) črez substituciju \(S\).} i črez
\[
(3)\qquad u_Su_T=u_{ST}a_{S,T}\quad \text{s }a_{S,T}\text{ v }K^\ast
\]
\[
\begin{gathered}
(4)\qquad a_{ST,R}\,a_{S,T}^{\,R}=a_{S,TR}\,a_{T,R}\\
\text{(zakon asociativnosti iz \([u_Su_T]u_R=u_S[u_Tu_R]\)).}
\end{gathered}
\]
\(A\) nazyvaje se skrěženy produkt od \(K\) s \(\mathfrak G\) pri faktornom sistemu \(a_{S,T}\); dokazuje se, že \(A\) stavaje prosta normalna algebra nad \(k\), to jest matrično kolco \(r\)-togo stepena \(D_r\) nad prirěđenoju dělitvenoju algebraju \(D\), i že \(K\) jest maksimalno komutativno podpolje, znači razpadno polje (t. j. razširjenje oblasti koeficientov \(k\) poljem izomorfnym k \(K\) daje razloženu algebru, polno matrično kolco nad centrom). Obratno, k zadanoj dělitvenoj algebrě \(D\) vsegda obstajajut matrične kolca \(D_r\), ktore možno porođati označenym sposobom kako skrěženy produkt.

Ako se prěhodi od \(u_S\) k \(v_S=u_Sc_S\) s \(c_S\) v \(K^\ast\), čto porođaje toj samy avtomorfizm, nastavajut ``asociovane'' faktorne sistemy:
\[
(5)\qquad \bar a_{S,T}=a_{S,T}\frac{c_S^{\,T}c_T}{c_{ST}}.
\]
Asociovane faktorne sistemy sobirajut se v jeden klas \((a)\); podobno sobirajut se vse algebry podobne k \(A\), t. j. vse \(D_r\) s \(r=1,2,\ldots\), v jeden klas \(\mathfrak A\). \emph{Klasy \(\mathfrak A\) i \((a)\) odpovědajut jedno-jednoznačno: klasy s fiksovanym razpadnym poljem \(K\) tvoret, vzhodno k direktnomu tvorjenju produkta, abelovu grupu, izomorfnu člen-po-členu produktu klasov faktornyh sistemov. Jediničny element jest klas razloženyh algeber odnosno sistem vsih veličin transformacije}
\[
        \frac{c_S^{\,T}c_T}{c_{ST}}.
\]
Ide tut o grupu klasov algeber, zaměčenu R. Brauerom.

\paragraph{4.}
Teper hču črez \emph{specializaciju} na \emph{ciklične razpadne polja} dojti k svęzi s \emph{ponjatjem normy}, i s tym k formulaciji obobčenogo \emph{normnogo teorema} po principu izloženom na početku. Ako \(Z\) jest ciklično, \(S\) porođujuča substitucija jego grupy -- odpovědajuča algebra togda nazyvaje se cikličnoju -- možno stepenjam od \(S\) dati odpovědajuče stepenje od \(u\), znači
\[
\begin{aligned}
(1')\qquad &A=Z+uZ+\cdots+u^{n-1}Z\\
(2')\qquad &zu=uz^S\\
(3')\quad &u^n=a\\
(4')\quad &a\text{ leži v osnovnom polju }k^\ast\\
(5')\quad &\bar a=a\cdot N(c),\quad \text{ako se postavi }v=uc.
\end{aligned}
\]
Vsaky faktorny sistem sostoji tut teda iz jednogo-jedinogo elementa \(a\), ležečego v osnovnom polju -- oznaka \(A=(a,Z)\). Jediničny klas faktornyh sistemov dan jest normami iz \(Z^\ast\), a grupa klasov algeber stavaje izomorfna faktor-grupě
\[
        k^\ast/N(Z^\ast).
\]
Ciklična algebra \((a,Z)\) razkladaje se togda i samo togda, kogda \(a\) jest norma někojego elementa iz \(Z\).

Ta svęz medžu normoju i razloženjem daje formulaciju ``normnogo teorema'', imenno \emph{teorem o razloženih algebrah: ako algebra razkladaje se v vsakom městě, togda ona razkladaje se bezuslovno}. Pri tom ``město'', kako obyčno v teoriji čisel, opreděljaje se tym, že osnovno polje \(k\) zaměnja se jego \(p\)-adičnym razširjenjem \(k_{\mathfrak p}\), gde \(\mathfrak p\) jest prosty ideal iz \(k\); odpovědajuče konečno mnogim beskonečnym městam, \(k\) i jego konjugovane razširjajut se poljem realnyh čisel.

Napravdu, v tom soderži se normny teorem za ciklične polja. Bo za ciklične algebry \((a,Z)\), po vyše rečenom, teorem jest ravnosilen tvrdženju: ako \(a\) v vsakom (konečnom ili beskonečnom) městě jest \(p\)-adična norma, togda \(a\) jest norma čisla iz \(Z\); ili, bez prěhoda k \(p\)-adičnomu, \emph{ako \(a\) jest normny ostatok po vsakom prostom idealu \(\mathfrak p\) iz \(k\) (i zadovolja někojim uslovjam znakov), togda \(a\) jest norma čisla iz \(Z\).} Ta poslednja forma jest právě normny teorem, dokazovany v teoriji klasovyh polj s upotrěboju znanyh analitičnyh sredstv. A dokaz občego teorema o razloženih algebrah možno dobiti iz togo cikličnogo specialnogo slučaja čisto algebraično-aritmetičnymi razsmotrenjami.\footnote{R. Brauer, H. Hasse, E. Noether, \foreign{\emph{Beweis eines Hauptsatzes in der Theorie der Algebren}}. J. f. M. 167 (1932).} Prvy važny poslědok iz toga vyvel Hasse: \emph{vsaka prosta normalna algebra nad algebraičnym čislovym poljem jest ciklična.} V poisku dokaza togo davno predpologanogo fakta nastala obča formulacija.

\paragraph{5.}
Drugy poslědok iz teorema o razloženih algebrah -- znovu čisto algebraično-aritmetičnymi zaključivanjami -- jest teorem o glavnom rodu, postavjeny na početku.\footnote{Dokaz ima iziti v \foreign{Math. Ann.}} Jego invariantna formulacija opiraje se na fakt, že relacije (2) do (5), ktore definišut skrěženy produkt, sut čisto multiplikativne i zato ostavajut smyslene, ako \(K^\ast\) zaměniti abelovoju grupoju \(\mathfrak J\), ktora zadovolja samo uslovju, že jej grupa avtomorfizmov soderži podgrupu izomorfnu k \(\mathfrak G\). Na město (1) togda nastupaje ``razširjenje \(\mathfrak G\) posrědstvom \(\mathfrak J\)'' v smyslu teorije grup. Ako za \(\mathfrak J\) vzęti grupu vsih idealov iz \(K\), togda faktorne sistemy stavajut sistemami idealov; razděljenje na klasy v \(\mathfrak J\) induciraje razděljenje na klasy za faktorne sistemy, i obće to jest tonše razděljenje na idealne klasy než izhodiščno. Bo trebovanje, aby množenje \(u_S\) s (absolutnymi) idealnymi klasami bylo jednoznačno, govori samo, že veličiny transformacije \(c_S^{\,T}c_T/c_{ST}\) -- jediničny klas elementnyh faktornyh sistemov -- ležet v jediničnom klasu idealnyh faktornyh sistemov. V dělě však, kako pokazuje specializacija k znanym slučajam, dosta jest uže něčto menje tonkogo razděljenja. Ja definišu: \emph{k jediničnomu klasu faktornyh sistemov pričisljuju se ti elementy \(a_{S,T}\), ktore vo vsih (konečnyh i beskonečnyh) městah razvětvjenja od \(K\) porođajut razložene algebry.}

Tako nastavajuče razširjenje od \(\mathfrak G\) bude označeno \(\mathfrak G^\ast\); \((c)\) označaje absolutny idealny klas od \(c\). Togda

\emph{invariantna formulacija teorema o glavnom rodu} glasi: \emph{ako pri tako opreděljenom razděljenju na klasy substitucija \(v_S=u_S(c_S)\) porođaje avtomorfizm od \(\mathfrak G^\ast\) -- vse te \((c_S)\) tvoret glavny rod --, togda toj avtomorfizm jest vnutrišny i porođaje se někojim idealnym klasom \((\mathfrak b)\).}

Znane specialne slučaje slědujut iz ravnosilnoj, něčto eksplicitnějšej formulacije: \emph{ako veličiny transformacije \((c_S^{\,T})(c_T)/(c_{ST})\), utvorjene iz idealnyh klasov \((c_S)\), naležet k jediničnomu idealnomu klasu faktornyh sistemov, togda obstaja idealny klas \((\mathfrak b)\) takovy, že \((c_S)\) stavajut symboličnymi \((1-S)\)-tymi stepenjami: \((c_S)=(\mathfrak b)/(\mathfrak b^{S})\) za vse \(S\) iz \(\mathfrak G\).} Bo predpostavjenje upravo izražaje svojstvo avtomorfizma; fakt, že on jest vnutrišny, izražaje se ravnostju \(v_S=(\mathfrak b)^{-1}u_S(\mathfrak b)=u_S(\mathfrak b)^{1-S}=u_S(c_S)\).

Specializacija k cikličnomu slučaju daje teda, s ohledom na normovanje: ako \(N([c])\) leži v jediničnom idealnom klasu faktornyh sistemov, togda \((c)\) stavaje symbolična \((1-S)\)-ta stepenj:
\[
(c)=(\mathfrak b)^{1-S}.
\]

\paragraph{6.}
Da by odtud dojti k znanoj formě za ciklične polja i kvadratne formy, najprvo zaměču, že teorem jest izrečen za polne idealne klasy, ale jego soderžanje ostava to samo, ako se, kako obyčno, ograničiti na idealy proste k městam razvětvjenja.

Togda tut opreděljeny glavny rod za kvadratne polja prěhodi v Gaussov. Napravdu, idealnym klasam odpovědajut kvadratne formy, a normam klasov čisla, ktore možno predstaviti tymi formami. To, že jediničny klas porođaje algebry razložene v městah razvětvjenja od \(K\), znači, že ta predstavlajema čisla v městah razvětvjenja stavajut kvadratnymi ostatkami; odpovědajuče formy zato imajut totalny karakter glavnoj formy, to jest tvoret Gaussov glavny rod. Znano jest, že symbolična \((1-S)\)-ta stepenj prěhodi v duplikaciju.

Za ciklične polja formulacija prěhodi v slědujuću: glavny rod sostoji iz vsih idealnyh klasov, ktoryh normy v (konečnyh i beskonečnyh) městah razvětvjenja stavajut normnymi ostatkami. To jest, kako znano, ravnosilno uslovju ``normny ostatok po provodniku'', i tym obyčny teorem nastavaje kako specializacija.

Mimo to možno takodže v občem slučaju libovoljnyh Galoisovyh polj vvesti provodnik, sostavjeny samo iz měst razvětvjenja, tak, že pri normovanju faktornyh sistemov za vsako iz tyh měst luč soderži samo elementy jediničnogo klasa.

Tut nastavaje vopros o svęzi s Artinovymi provodnikami, spomenutymi v pregledu 2, bo oni sostavjajut se iz tyh samyh prostyh idealov; a vměstě s tym -- vopros o svęzi s teorijeju Galoisovyh modulov, drugoju hiperkompleksnoju metodoju. Kako daleko budut segati te dvě metody, pokaže tek budučnost.

\end{document}
